% ===========================================================================
%  sviluppatore/05-estendere/02-strategia.tex
% ===========================================================================
\section{Una strategia}\label{sez:nuova-strategia}

\subsection{Che cosa si crea}

\begin{codiceInterno}{console}{La struttura minima}{strategies/}
strategies/mia_strategia/
|-- __init__.py       # ri-esporta DESCRIPTOR
|-- strategy.py       # espone DESCRIPTOR
`-- mia_impl.py       # l'implementazione
\end{codiceInterno}

\noindent
La cartella \textbf{non} deve cominciare con \cod{\_}: il prefisso è
l'interruttore di spegnimento, e il filtro che lo applica è questo.

\begin{codiceInterno}{python}{Il filtro delle cartelle}{strategies/registry.py:55}
def _is_valid_strategy_dir(p: Path) -> bool:
    if not p.is_dir():                                   return False
    if p.name.startswith("_") or p.name.startswith("."): return False
    if p.name in SKIP_DIR_NAMES:                         return False
    if not (p / "strategy.py").exists():                 return False
    return True
\end{codiceInterno}

\subsection{Il livello di esecuzione}

Si sottoclassa \simbolo{BaseAtomicStrategy} da \file{strategies/base.py}, che
offre già \simbolo{applicable\_symbols()}, \simbolo{get\_value()} e
\simbolo{get\_last\_value\_until()}. Tre cose sono obbligatorie:

\begin{codiceInterno}{python}{Il contratto di esecuzione}{strategies/base.py:55}
class AtomicStrategy(ABC):
    @property
    @abstractmethod
    def style(self) -> StrategyStyle: ...

    @property
    @abstractmethod
    def allowed_asset_classes(self) -> set[AssetClass]: ...

    @abstractmethod
    def generate_signals(
        self,
        t_decision: datetime,
        t_execution: datetime,
        indicators_by_symbol: dict[str, PrecomputedIndicators],
        portfolio: Portfolio,
    ) -> list[Signal]:
\end{codiceInterno}

\noindent
I valori ammessi per \cod{style} vengono da \simbolo{core.types.StrategyStyle}:
\cod{TREND}, \cod{MEAN\_REVERSION}, \cod{MOMENTUM}, \cod{BREAKOUT},
\cod{STAT\_ARB}, \cod{STAT\_ARB\_INTRADAY}, \cod{FACTOR}, \cod{VALUE},
\cod{MACRO}, \cod{CARRY}, \cod{VOL}.

\begin{attenzione}{I nomi dei parametri fanno parte del contratto}
Il motore chiama \cod{generate\_signals} per \emph{keyword}
(\vediSez{decisione-esecuzione}). Rinominare \cod{t\_decision} in \cod{t}, o
\cod{indicators\_by\_symbol} in \cod{ind}, fa fermare il backtest alla prima
barra con un \cod{TypeError}.
\end{attenzione}

\subsection{Il livello dei metadati}

\simbolo{StrategyDescriptor} ha sei campi obbligatori e una ventina di
opzionali.

\begin{codiceInterno}{python}{I sei obbligatori}{strategies/descriptor.py:107}
    # Required
    name: str
    description: str
    category: StrategyCategory
    timeframe: StrategyTimeframe
    factory: Callable[[dict], Any]
    default_params: dict
\end{codiceInterno}

\noindent
Un esempio reale, che mostra anche il modo idiomatico di scrivere la
\cod{factory}:

\begin{codiceInterno}{python}{Il descrittore di una strategia installata}{strategies/momentum/strategy.py:16}
DESCRIPTOR = StrategyDescriptor(
    name="momentum",
    description="Cross-Sectional Momentum 12-1 (Jegadeesh-Titman). "
                "Top decile, monthly rebalance, long-only.",
    category=StrategyCategory.D1_LONG_ONLY,
    timeframe=StrategyTimeframe.D1,

    excluded_instruments=["EURUSD", "GBPUSD", "USDJPY", "USDCHF", "AUDUSD"],
    required_indicators=["momentum_252_21"],

    factory=lambda params: CrossSectionalMomentumStrategy(**{
        k: v for k, v in params.items() if not k.startswith("exit_")
    }),
    default_params={
        "formation_window": 252,
        "skip_window": 21,
        "top_decile_pct": 0.10,
        "max_positions": 25,
        "rebalance_every_n_days": 21,
        "exit_catastrophic_sl_pct": None,
    },

    min_lookback_days=300,
    backtest_script="strategies/momentum/run_backtest.py",
    output_path_pattern="output/strategies/wf_momentum_oos_equity*.parquet",
    stream_type="equity_curve",
    stream_category="d1",
    exit_policy_type="signal_only",
    regime_weights={"bull": 1.0, "neutral": 1.0, "bear": 0.7},
)
\end{codiceInterno}

\begin{nota}{Il filtro delle chiavi \cod{exit\_}}
La \cod{factory} scarta le chiavi che cominciano con \cod{exit\_} prima di
passarle al costruttore. Quei parametri non appartengono alla strategia: li
consuma la politica di uscita, costruita a parte da
\cod{exit\_policy\_type}. Scrivere la \cod{factory} senza quel filtro
produce un \cod{TypeError} su un argomento inatteso.
\end{nota}

\subsection{I campi opzionali che contano davvero}

\begin{center}
\begin{tabular}{@{}L{40mm}L{98mm}@{}}
\toprule
\textbf{Campo} & \textbf{A che serve} \\
\midrule
\rigaTenue
\cod{output\_path\_pattern},\newline\cod{stream\_type},
\cod{stream\_category}
  & \textbf{Integrazione con la dashboard, senza toccarla.} Con questi tre
    campi lo stream compare da solo in \voce{Streams}: non c'è nessun
    elenco da aggiornare in \file{dashboard/config.py}. Senza
    \cod{output\_path\_pattern} lo stream non compare. \\
\cod{required\_indicators}
  & Nomi presi dal pool comune. Il registro ne calcola l'unione su tutte le
    strategie e li pre-calcola una volta sola. \\
\rigaTenue
\cod{regime\_weights}
  & Modulatori \cod{bull}/\cod{neutral}/\cod{bear}. Le chiavi devono essere
    esattamente quelle tre e i valori numerici non negativi: è l'unico
    controllo del descrittore che solleva \cod{ValueError} invece di un
    avviso. \\
\cod{allocator\_overrides}
  & Scavalca per questa sola strategia \cod{max\_position\_weight},
    \cod{max\_gross\_leverage}, \cod{max\_asset\_class\_gross},
    \cod{portfolio\_target\_vol\_annual}. \\
\rigaTenue
\cod{extra\_timeframes}
  & Per le strategie multi-timeframe. Il caricamento parte dal più fine. \\
\cod{required\_external\_data}
  & \cod{news}, \cod{fundamentals}, \cod{policy\_rates}. Il sistema li
    \emph{scarica} se mancano. \\
\rigaTenue
\cod{required\_mms\_features},\newline\cod{required\_volume\_profile}
  & Feature specialistiche calcolate su richiesta. \\
\cod{maintenance}
  & \cod{retrain\_every\_days}, \cod{retrain\_every\_bars},
    \cod{retrain\_on\_drawdown\_pct}, \cod{data\_refresh\_every\_days}. È
    da qui che il sistema sa quando un modello va riaddestrato. \\
\rigaTenue
\cod{\_module\_path}, \cod{\_folder\_path}
  & \textbf{Li scrive il registro.} Non vanno impostati a mano. \\
\bottomrule
\end{tabular}
\end{center}

\begin{attenzione}{\cod{regime\_weights} si dichiara a priori}
Il commento nel descrittore è esplicito: modificarli guardando il backtest
«sarebbe curve-fitting mascherato da regime adaptation». Si dichiarano sul
ragionamento strutturale --- questa strategia compra il momentum, quindi in
regime ribassista deve pesare meno --- e poi si verifica che il modulato non
peggiori rispetto al neutro. Non si cercano i valori che massimizzano il
risultato.
\end{attenzione}

\subsection{La procedura}

\begin{pratica}{Aggiungere una strategia}
\begin{enumerate}
  \item Creare la cartella \textbf{senza} prefisso \cod{\_}.
  \item \file{mia\_impl.py}: sottoclassare \simbolo{BaseAtomicStrategy},
        implementare \cod{style}, \cod{allowed\_asset\_classes},
        \cod{reset()} e \cod{generate\_signals()} \emph{con quei nomi esatti
        di parametro}.
  \item Leggere solo fino a \cod{t\_decision}. Mettere
        \cod{Signal.timestamp = t\_execution}.
  \item \file{strategy.py}: \cod{DESCRIPTOR = StrategyDescriptor(...)} con i
        sei campi obbligatori; \cod{factory} che filtra le chiavi
        \cod{exit\_}.
  \item Per vedere lo stream in dashboard: \cod{output\_path\_pattern} +
        \cod{stream\_type} + \cod{stream\_category}.
  \item \file{\_\_init\_\_.py}: ri-esportare \cod{DESCRIPTOR}.
  \item Verificare con \file{tests/diagnostica\_strategie.py}, che esiste
        proprio per rispondere alla domanda «perché
        \cod{discover\_strategies()} torna vuota».
\end{enumerate}
\end{pratica}

\noindent
Il riferimento completo è \file{strategies/GUIDE\_STRATEGIE\_CUSTOM.md}
(\num{2281} righe), affiancato da sei guide tematiche: il grafo a nodi come
alternativa al Python, l'uso dei regimi dentro una strategia, le posizioni
multiple sullo stesso simbolo, le feature MMS, il volume profile e un
documento di proposta sulle strategie in valuta.
